\scanpage{64}
Les morphismes les plus importants de la
situation apparaissent dans le diagramme
\inserted{hexagonal} suivant
\diagram{\begin{tikzcd}[column sep=small,row sep=large,execute at end picture={
\draw[densely dotted] ([xshift=7pt]s64upper.east) arc[start angle=270,end angle=90,radius=3pt];
\draw[densely dotted] ([xshift=7pt]s64centre.east) arc[start angle=270,end angle=90,radius=3pt];
\draw[densely dotted] ([xshift=-5pt,yshift=5pt]s64right.north west) arc[start angle=135,end angle=315,radius=3pt];
\draw[densely dotted] ([xshift=7pt]s64lower.east) arc[start angle=270,end angle=90,radius=3pt];}]
&|[alias=s64upper]|\Sigma^{<}(1)\arrow[rr,"\uline{b}",start anchor=east,shorten <=7pt]\arrow[dl,"\uline{s}"']&&\Sigma&\\
\Sigma&&|[alias=s64centre]|\Sigma^{<}(2)\arrow[ul,hook,"\uline{s'}"]\arrow[rr,"\uline{b'}",start anchor=east,shorten <=7pt]\arrow[dl,"p"']\arrow[ur,phantom,"\text{cart}" description]\arrow[dr,phantom,"\text{cart}" description]&&|[alias=s64right]|\Sigma^{<}(1)\arrow[ul,"\uline{b}"',start anchor=north west,shorten <=7pt]\arrow[dl,"\uline{s}"]\\
&|[alias=s64lower]|\Sigma^{<}(1)\arrow[ul,"\uline{s}"]\arrow[rr,"\uline{b}"',start anchor=east,shorten <=7pt]&&\Sigma&
\end{tikzcd}},
(40)

où les flèches marquées \diagram{\begin{tikzcd}[column sep=large,execute at end picture={
\draw[densely dotted] ([xshift=7pt]s64legend.east) arc[start angle=270,end angle=90,radius=3pt];}]
|[alias=s64legend]|{}\arrow[r,start anchor=east,shorten <=7pt]&{}
\end{tikzcd}} (en fait,
$\uline{b},\deleted{\unsure{p_2}},\uline{b'},\uline{s'}$) sont des immersions bordantes, \inserted{propres}
et les flèches marquées $\longrightarrow$ (en fait,
$\uline{s}$ et $p$) sont (dans le cas équisingulier)
fibrantes et propres \unsure{\}} (voire, de plus : fibres
des variétés à bord, dont le bord est
dans le cas de $\uline{s}:\Sigma^{<}(1)\longrightarrow\Sigma$, l'image
de $\Sigma^{<}(2)\overset{\uline{s'}}{-}\Sigma^{<}(1)$\inserted{(i.e. l'ens des $(x\prec y)$ tels que $y$ \uline{ne soit pas} successeur de $x$)}).
\margin{donc, dans le cas de $p:\Sigma^{<}(2)\longrightarrow\Sigma^{<}(1)$,
le fibré des bords est l'image de $\Sigma^{<}(3)$}
On fera attention
que les applications fibrantes $s,p$ n'ont \deleted{\unsure{rien}}\inserted{rien}
de surjectif --- l'image de $s$ est formé des $x\in\Sigma$
qui \uline{ne sont pas maximaux}, celle de $p$ est formé
des $(x\prec y)$ tels que $y$ ne soit pas maximal.
\mspage{63}
